% Part: incompleteness
% Chapter: introduction
% Section: decidability

\documentclass[../../../include/open-logic-section]{subfiles}

\begin{document}

\olfileid[fr]{inc}{int}{dec}

\olsection{Indécidabilité et incomplétude}

La démonstration des théorèmes d'incomplétude de Gödel exige une
arithmétisation de la syntaxe. Pourtant, même sans celle-ci, nous
pouvons déjà obtenir des résultats intéressants en supposant seulement
qu'une théorie !!{represents} toutes les relations !!{decidable}s.
Cette preuve repose sur un argument diagonal analogue à celui qui
établit l'indécidabilité du problème de l'arrêt.

\begin{thm}
Si $\Gamma$ est une théorie cohérente qui !!{represents} toute
relation !!{decidable}, alors $\Gamma$ n'est pas !!{decidable}.
\end{thm}

\begin{proof}
Supposons $\Gamma$ !!{decidable}. Nous allons montrer que si
$\Gamma$ !!{represents} toute relation !!{decidable}, elle est
incohérente.

Les propriétés décidables (relations unaires) sont représentées
par des !!{formula}s ayant pour seule variable libre~$x$. Soit
$!A_0(x)$, $!A_1(x)$, \dots, une énumération calculable de toutes
les !!{formula}s de ce langage ayant cette variable libre, qu'elles représentent ou non une propriété !!{decidable}.\footnote{Précision éditoriale : l'énumération porte sur la syntaxe des formules, dont la reconnaissance est effective, et non sur la classe sémantique des seules représentations de propriétés décidables.} Considérons alors l'ensemble $D\subseteq\Nat$ :
\[
D = \Setabs{n}{\Gamma \Proves \lnot !A_n(\num{n})}
\]
$D$ est !!{decidable} : pour savoir si $n\in D$, nous calculons
d'abord $!A_n(x)$, puis $\lnot !A_n(\num{n})$. La substitution du
terme $\num{n}$ à chaque occurrence libre de~$x$ dans $!A_n(x)$,
puis l'ajout de $\lnot$ devant $!A_n(\num{n})$, sont des opérations
mécaniques.\footnote{Correction éditoriale : dans deux occurrences du calcul source, $!A$ doit rester $!A_n$ ; l'indice est conservé dans toute la substitution et le test d'appartenance.} Par hypothèse, $\Gamma$ est !!{decidable} : nous pouvons
donc déterminer si $\lnot !A_n(\num{n})\in\Gamma$. Si oui,
$n\in D$; sinon, $n\notin D$. Ainsi, $D$ est !!{decidable}.

Puisque $\Gamma$ !!{represents} toutes les propriétés !!{decidable}s,
elle !!{represents}~$D$. Toute !!{formula} qui représente~$D$
dans~$\Gamma$ figure parmi $!A_0(x)$, $!A_1(x)$, \dots.
Choisissons donc $d$ tel que $!A_d(x)$ représente~$D$ dans~$\Gamma$.
Si $d\notin D$, la représentation de~$D$ donne
$\Gamma\Proves\lnot !A_d(\num{d})$. Or cela signifie, par la définition
de~$D$, que $d\in D$, contrairement à notre hypothèse. Donc
$d\in D$.

Puisque $d\in D$, la définition de~$D$ donne
$\Gamma\Proves\lnot !A_d(\num{d})$. D'autre part, comme
$!A_d(x)$ représente~$D$ dans~$\Gamma$, nous avons
$\Gamma\Proves !A_d(\num{d})$. Ainsi, $\Gamma$ est incohérente.
\end{proof}

\begin{explain}
Ce théorème montre qu'aucune théorie cohérente représentant toutes
les relations !!{decidable}s ne peut être !!{decidable}. Nous
montrerons que $\Th{Q}$ !!{represents} effectivement toutes ces
relations ; toute extension de~$\Th{Q}$ les représente donc aussi.
Toute extension \emph{cohérente} de~$\Th{Q}$ est ainsi indécidable.\footnote{Précision éditoriale : la source ne précise pas ici la cohérence et semble attribuer un modèle standard à toutes les extensions. La vérité dans le modèle standard concerne les trois théories citées ; une extension arbitraire de~$\Th{Q}$ peut être incohérente et donc décidable.}
En particulier, tous les énoncés de $\Th{Q}$, $\Th{PA}$ et $\Th{TA}$ sont vrais
dans le modèle standard ; ces trois théories sont donc cohérentes et indécidables.

Nous pouvons aussi en déduire une forme faible du premier théorème
d'incomplétude. Toute théorie !!{axiomatizable} et complète est
!!{decidable}. Une théorie cohérente qui est !!{axiomatizable} et
représente toutes les propriétés !!{decidable}s ne peut donc pas être
complète.
\end{explain}

\begin{thm}
Si $\Gamma$ est !!{axiomatizable} et complète, elle est
!!{decidable}.
\end{thm}

\begin{proof}
Toute théorie incohérente est !!{decidable}, puisqu'elle contient
tous les !!{sentence}s : à la question « $!A\in\Gamma$ ? », on
répond toujours « oui ».

Supposons donc $\Gamma$ cohérente, !!{axiomatizable} et complète.
Comme elle est !!{axiomatizable}, l'ensemble de ses théorèmes est
!!{computably enumerable}. En effet, une procédure calculable peut
énumérer toutes les !!{derivation}s correctes à partir des axiomes
de~$\Gamma$. De chacune, on peut calculer l'!!{sentence} dérivé;
on obtient ainsi une énumération calculable des théorèmes de~$\Gamma$.
Puisque $\Gamma$ est cohérente et complète, $!A$ est un théorème
si et seulement si $\lnot !A$ n'en est pas un. Pour décider si
$!A\in\Gamma$, énumérons alors les théorèmes de~$\Gamma$.
Si $!A$ apparaît, nous savons que $\Gamma\Proves !A$;
si $\lnot !A$ apparaît, nous savons que $\Gamma\Proves/ !A$.
L'un de ces deux cas finit nécessairement par se produire, puisque
$\Gamma$ est complète. La procédure donne donc toujours une réponse.
\end{proof}

\begin{cor}
\ollabel{cor:incompleteness}
Si $\Gamma$ est cohérente, !!{axiomatizable} et représente toute
propriété !!{decidable}, elle n'est pas complète.
\end{cor}

\begin{proof}
Si $\Gamma$ était complète, le théorème précédent montrerait qu'elle
est !!{decidable}, puisqu'elle est !!{axiomatizable} et cohérente.
Mais le premier théorème montre qu'une théorie qui représente toute
propriété !!{decidable} n'est pas !!{decidable}. Contradiction.
\end{proof}

\begin{prob}
Montrez que $\Th{TA}=\Setabs{!A}{\Sat{N}{!A}}$ n'est pas
!!{axiomatizable}. Vous pouvez admettre que $\Th{TA}$ représente
toutes les propriétés décidables.
\end{prob}

Une fois établi que $\Th{Q}$ représente toutes les propriétés
!!{decidable}s, le corollaire montrera que $\Th{Q}$ est incomplète.
Cependant, sa preuve ne donne aucun exemple d'!!{sentence}
indépendant : elle montre seulement qu'il en existe. Pour en construire
un sous la seule cohérence, il faudra arithmétiser la syntaxe et utiliser
la construction de Rosser, issue de l'idée originale de Gödel. Il restera aussi à démontrer que $\Th{Q}$ représente
effectivement toutes les propriétés !!{decidable}s.

Toute théorie mathématiquement \emph{intéressante} n'est cependant pas
nécessairement incomplète ou indécidable. Nombre de théories assez
riches pour décrire des faits mathématiques intéressants ne satisfont
pas les hypothèses du résultat de Gödel. Par exemple, la théorie
$\Th{Pres}=\Setabs{!A\in\Lang{L_{A^+}}}{\Sat{N}{!A}}$, constituée
des !!{sentence}s du langage de l'arithmétique sans~$\times$ qui sont
vrais dans le modèle standard, est à la fois complète et !!{decidable}.
C'est l'\emph{arithmétique de Presburger}; elle démontre toutes les
vérités sur les nombres naturels exprimables avec les seuls symboles
$\Obj{0}$, $\prime$ et~$+$. Le symbole~$<$ qui figure dans~$\Lang{L_{A^+}}$ est définissable à partir de ces symboles.

\end{document}
